% Split from 7-1EleanorMia.tex by cut-sheet v2/v3 — content below is the source scene, header adjusted
\chapter{Eleanor's Writing Plea II: The Late-Night Chessboard}
% \section*{A Late-Night Chessboard Problem}

After Mia had retired to bed, her notebook riddled with childlike sketches of pyramids and cubes, Eleanor remained by the fireplace. Flames a dancing across the room, her mind elsewhere, entangled in webs of misplaced number patterns. Then a faint gleeful sound from the corner, no more a burst of modest jubilation. She turned to see Mo, Mia's older brother sitting cross-legged with his laptop, a rare smile lighting up his face. Eleanor walked over, half-expecting to find him skulking iver some age inappropriate , she saw the ChessKid interface on his screen. "Checkmate," he exclaimed, leaning back triumphantly.

Relieved, Eleanor smiled. "Well done, Mo! Looks like you’re quite the strategist. But online chess is easy can you handle looking your opponent in the eye?" Mo shrugged, still grinning. "A win’s a win, Auntie Eleanor."

"True," she said, glancing at the nearby bookshelf. An idea struck her. "How about we go analog for a change? Let’s see how you fare on a real chessboard."
Mo’s face lit up. "You’re on!" he said, springing up to help her fetch the old wooden chess set. 

"Mo," she began, her voice taking on that tone of inquiry that made Mia roll her eyes, "have you ever thought about how many squares are on a chessboard?"

\begin{figure}[H]
\includegraphics[width=\linewidth]{Illustrations/vign-Mo1.png}
\caption{Mo after just getting his head around a 2-d problem}
\end{figure}


Mo looked up, puzzled. "Isn’t it just \(8 \times 8 = 64\)?"

Eleanor smiled mischievously. "Sure, for the smallest squares. But what about larger ones? What about \(2 \times 2\), or \(3 \times 3\)?"

Mo paused. "I… I don’t know. How would you even figure that out?"


\subsection*{Counting the Squares}

Eleanor reached for her notebook. "Let’s think about it step by step," she said, sketching a simple \(8 \times 8\) chessboard. "For \(1 \times 1\) squares, you’re absolutely right: there are 64 of them. But for \(2 \times 2\) squares, you can only fit them in a \(7 \times 7\) grid because each square takes up more space."

Mo nodded, starting to catch on. "So for \(3 \times 3\), it would be \(6 \times 6\) squares?"

"Exactly!" Eleanor said, her enthusiasm growing. "In general, for \(k \times k\) squares, you can fit \((8 - k + 1) \times (8 - k + 1)\). The total number of squares is the sum of all these possibilities."

She wrote out the formula:
\[
\text{Total squares} = \sum_{k=1}^{8} (8 - k + 1)^2.
\]

Mo leaned over, watching intently as she expanded the terms:
\[
\text{Total squares} = 8^2 + 7^2 + 6^2 + \ldots + 1^2.
\]

He frowned. "That’s a lot of adding."

Eleanor chuckled. "It is, but there’s a shortcut. Remember Mo the formula for the sum of squares we discussed the other day?". Here it is: 
\[
\sum_{k=1}^n k^2 = \frac{n(n+1)(2n+1)}{6}.
\]
"Remember?", Mo says "Sure" with some convincing assiduity:
"For \(n = 8\), we get:
\[
\text{Total squares} = \frac{8 \cdot 9 \cdot 17}{6} = 204.
\]
\noindent
 Eleanor thus completed the picture.
\begin{figure}[H]\centering
\includegraphics[width=0.7\textwidth]{Illustrations/chessboard.png}
\caption{How many squares in a \(8\times 8\) chessboard.}
\label {Chessboard:fig}
\end{figure}
The correct total number of squares on an 8x8 chessboard, including all sizes of squares from 1x1 to 8x8, is calculated as the sum of the first eight square numbers:
\begin{align*}
\text{For 1x1 squares: } & 8 \times 8 = 64 \\
\text{For 2x2 squares: } & 7 \times 7 = 49 \\
\text{For 3x3 squares: } & 6 \times 6 = 36 \\
\text{For 4x4 squares: } & 5 \times 5 = 25 \\
\text{For 8x8 squares: } & etc
\end{align*}
\text{Adding these gives: } 64 + 49 + 36 + 25 + 16 + 9 + 4 + 1 = 204. 
Mo’s eyes widened. "Wow, so there are 204 squares on a chessboard, not just 64?"

"Exactly," Eleanor pops, "It’s all about perspective. The problem looks simple at first, but when you dig a little deeper, there’s so much more to uncover."

% >>>>>> ANNEX A (cut-sheet v2): the hyper-dimensional pyramid extension — block commented out below, pointer left in text <<<<<<
\textit{(Where the pyramids go when the dimensions keep climbing is a game for the annexe --- Annexe A --- not for bedtime.)}

% \section*{Eleanor going hyper}
% 
% Eleanor paused as she is disinclined to do, watching Mo’s triumphant expression as he proceeded to try count the 204 squares aloud. Trying not to let his triteness dampen her teaching moment she reached for a Rubik’s Cube resting on the bookshelf and slapped it down in the middle of the chessboard, displacing pawns and rooks alike.
% 
% "Alright, Mo," she said with a grin, "now let’s take this a step further. What if I asked you how many cubes are in a \(8 \times 8 \times 8\) Rubik’s Cube?"
% 
% Mo stared at the colorful object, his earlier confidence waning. "You mean, like all the little cubes inside? That sounds… complicated."
% 
% Eleanor chuckled, realizing she might be pushing it. "It is. But let me share a story instead. You see, there’s a fascinating lesser-known cousin to Pascal’s Triangle—called Worpitzky’s Triangle, or as I like to call it, the ‘Power Pascal Triangle.’" Sketched the triangle as she spoke:
% \[
% \begin{array}{ccccccccccc}
% & & & & & 1_1 & & & & &\\
% & & & & 1_1 & & 1_1 & & & &\\
% & & & 1_1 & & 3_2 & & 2_3 & & &\\
% & & 1_1 & & 7_2 & & 12_3 & & 6_4 & &\\
% & 1_1 & & 15_2 & & 50_3 & & 60_4 & & 24_5 &\\
% 1_1 & & 31_2 & & 180_3 & & 390_4 & & 360_5 & & 120_6 \\
% & & \vdots & & & \vdots & & & \vdots & & \\
% \end{array}
% \]
% 
% "See how each number is generated?" she asked. "You take the two numbers directly above it, multiply them by their respective column positions, and sum them. Like this:
% \[
% 7_2 + 12_3 = 7 \times 2 + 12 \times 3 = 50_3.
% \]
% \newpage
% \noindent
% Now just like Pascal numbers 
% 
% \begin{minipage}[t]{0.5\textwidth}
% \[
% \begin{array}{ccccccccccc}
% & & & & & \binom{0}{0} = 1 & & & & & \\
% & & & & \binom{1}{0} = 1 & & \binom{1}{1} = 1 & & & & \\
% & & & \binom{2}{0} = 1 & & \binom{2}{1} = 2 & & \binom{2}{2} = 1 & & & \\
% & & \binom{3}{0} = 1 & & \binom{3}{1} = 3 & & \binom{3}{2} = 3 & & \binom{3}{3} = 1 & & \\
% & & \vdots & & & \vdots& & & \vdots & & \\
% \end{array}
% \]
% \end{minipage}
% 
% \noindent
% we can label the Worpitzky coefficients as:
% 
% \begin{minipage}[t]{0.5\textwidth}
% \[
% \begin{array}{ccccccccccc}
% & & & & & \worp{0}{0} & & & & & \\
% & & & & \worp{1}{0} & & \worp{1}{1} & & & & \\
% & & & \worp{2}{0} & & \worp{2}{1} & & \worp{2}{2} & & & \\
% & & \worp{3}{0} & & \worp{3}{1} & & \worp{3}{2} & & \worp{3}{3} & & \\
% & & \vdots & & & \vdots & & & \vdots & & \\
% \end{array}
% \]
% \end{minipage}
% 
% These numbers help us compute the sum of powers up to \(n^k\). 
% For example, for the sum of squares, we use the third row of the triangle:
% \[
% 1^2 + 2^2 + 3^2 + \dots + n^2 = \worp{2}{0} \binom{n}{1} + \worp{2}{1} \binom{n}{2} + \worp{2}{2} \binom{n}{3}.
% \]
% 
% Mo leaned forward, intrigued. "So you’re saying this triangle can tell us the total number of cubes?"
% 
% "Exactly!" Eleanor said, her enthusiasm unchecked. 
% 
% "For a \(8 \times 8 \times 8\) cube, the total number of cubes is just the sum of cubes up to \(8^3\). Let me show you how it works." Eleanor continued:
% \[
% 1^0 + 2^0 + \dots + 8^0 = 1 \binom{8}{1} = 8,
% \]
% \[
% 1^1 + 2^1 + \dots + 8^1 = 1 \binom{8}{1} + 1 \binom{8}{2} = 8 + 28 = 36,
% \]
% \[
% 1^2 + 2^2 + \dots + 8^2 = 1 \binom{8}{1} + 3 \binom{8}{2} + 2 \binom{8}{3} = 8 + 84 + 112 = 204,
% \]
% There is our number of squares in the chessboard. We have then for the cube:
% \begin{align*} 
% 1^3 + 2^3 + \dots + 8^3 &= 1 \binom{8}{1} + 7 \binom{8}{2} + 12 \binom{8}{3} + 6 \binom{8}{4} \\ &= 8 + 196 + 672 + 420 = 1296.
% \end{align*}
% Mo stared at the numbers, his eyes wide. "Wait… so for a Rubik’s Cube, you’re telling me there are 1296 smaller cubes?"
% 
% Eleanor shook her head, smiling gently. "Not quite, Mo. A Rubik’s Cube is a \(3 \times 3 \times 3\) cube, not \(8 \times 8 \times 8\). The number of smaller cubes is much simpler to calculate. It’s the sum of cubes up to \(n = 3\):"
% 
% \[
% 1^3 + 2^3 + 3^3 = \worp{3}{0} \binom{3}{1} + \worp{3}{1} \binom{3}{2} + \worp{3}{2} \binom{3}{3} + \worp{3}{3} \binom{3}{4}.
% \]
% 
% Eleanor pointed to the coefficients in the Worpitzky triangle:
% \[
% \worp{3}{0} = 1, \quad \worp{3}{1} = 7, \quad \worp{3}{2} = 12, \quad \worp{3}{3} = 6.
% \]
% 
% Plugging in the binomial coefficients:
% \begin{align*} 
% 1 \cdot \binom{3}{1} + 7 \cdot \binom{3}{2} + 12 \cdot \binom{3}{3} + 6 \cdot \binom{3}{4}&=1 \cdot 3 + 7 \cdot 3 + 12 \cdot 1 + 6 \cdot 0 \\&= 3 + 21 + 12 = 36.
% \end{align*}
% "So," Eleanor intent on wrapping this up now, as she had some adult thoughts of her own to tend to, "a \(3 \times 3 \times 3\) Rubik’s Cube has exactly \(36\) smaller cubes. "The power of recursion and Pascal-like structures at work. And this is just the beginning. Worpitzky’s Triangle connects deeply to sums of higher powers and even hints at a recursive beauty between them."
% 
% Mo scratched at his raggedy arm, overwhelmed a little but knowing enough to show he was impressed. "Auntie Eleanor, you’ve officially blown my mind."
% 
% She laughed, ruffling his hair. "That’s what mathematicians do, Mo. We find joy in puzzles, patterns, and just a little chaos."
% 
% Mo grinned, clearly inspired. "Okay, Auntie Eleanor. Your turn. My teacher asked us the other day, how many ways are there for a rook or a king to move from one corner of the board to the far corner?. She suggested I try to count the ways systematically and spot a pattern."
% 
% >>>>>> end of annexed block <<<<<<
% >>>>>> ANNEX A (cut-sheet v2): the King's Walk combinatorics — block commented out below, pointer left in text <<<<<<
\textit{(The King's Walk combinatorics --- the full working is set out in Annexe A.)}

% \section*{A King's Walk}
%  
% Mo's question hung in the air like a small dare. Eleanor considered him for a moment, because there are children to whom you give the simple answer and children to whom you give the answer that will haunt them, and she had long since worked out which one Mo was.
% \begin{figure}[h]
% \centering
% \begin{tikzpicture}[scale=0.55]
%   % Chessboard squares
%   \foreach \i in {0,...,7} {
%     \foreach \j in {0,...,7} {
%       \pgfmathparse{int(mod(\i+\j,2))}
%       \ifnum\pgfmathresult=0
%         \fill[boxbg!60] (\i,\j) rectangle (\i+1,\j+1);
%       \else
%         \fill[white] (\i,\j) rectangle (\i+1,\j+1);
%       \fi
%     }
%   }
%   % Border
%   \draw[gold, line width=0.7pt] (0,0) rectangle (8,8);
%   % Axis tick labels
%   \foreach \i/\label in {0/a,1/b,2/c,3/d,4/e,5/f,6/g,7/h} {
%     \node[below, font=\scriptsize\itshape, color=midgrey] at (\i+0.5,-0.05) {\label};
%   }
%   \foreach \j/\label in {0/1,1/2,2/3,3/4,4/5,5/6,6/7,7/8} {
%     \node[left, font=\scriptsize\itshape, color=midgrey] at (-0.05,\j+0.5) {\label};
%   }
% 
%   % King at origin
%   \fill[gold!90!black] (0.5,0.5) circle (0.32);
%   \node[font=\small\bfseries, white] at (0.5,0.5) {K};
%   \node[below right, font=\footnotesize\itshape, color=midgrey] at (-1.0,0.15) {$(0,0)$};
% 
%   % Three legal moves from (0,0)
%   \draw[->, forward, line width=1.3pt, >=Stealth]
%     (0.5,0.5) -- (1.45,0.5) node[midway, below, font=\scriptsize\itshape, color=forward] {E};
%   \draw[->, forward, line width=1.3pt, >=Stealth]
%     (0.5,0.5) -- (0.5,1.45) node[midway, left, font=\scriptsize\itshape, color=forward] {N};
%   \draw[->, joinclr, line width=1.3pt, >=Stealth]
%     (0.5,0.5) -- (1.45,1.45) node[midway, above left, font=\scriptsize\itshape, color=joinclr] {NE};
% 
%   % Destination
%   \draw[joinclr, line width=0.8pt, dashed] (7.5,7.5) circle (0.36);
%   \node[font=\small\bfseries, color=joinclr] at (7.5,7.5) {$\star$};
%   \node[above left, font=\footnotesize\itshape, color=midgrey] at (8.7,7.7) {$(n,n)$};
% 
%   % Geodesic dashed line
%   \draw[joinclr!60, line width=0.4pt, dashed] (0.5,0.5) -- (7.5,7.5);
% 
%   % An example path
%     \draw[<-, line width=1.0pt, opacity=0.7]
%     (0.5,0.5) -- (1.5,1.5) -- (2.5,1.5) -- (3.5,2.5) -- (3.5,3.5)
%     -- (4.5,4.5) -- (5.5,4.5) -- (5.5,5.5) -- (6.5,6.5) -- (7.5,7.5);
% \end{tikzpicture}
% \caption{The king walks from a1 to h8 using only East, North, and North-East
% moves. Three legal first moves are shown in colour; one example completed
% path is drawn faintly across the board. The dashed line marks the Euclidean
% geodesic.}
% \label{fig:chessboard}
% \end{figure}
%  
% ``A rook is easy,'' she said. ``It goes in straight lines. To get
% from a1 to h8 the rook must travel seven squares east and seven
% squares north, in some order. So the \href{https://leelemacjames.github.io/Dellanoy_Intro.pdf}{question} is just: how many ways
% can you arrange seven Es and seven Ns in a row of fourteen letters?''
%  
% Mo's lips moved. ``Fourteen choose seven.''
%  
% ``Three thousand four hundred and thirty-two,'' Eleanor said.
% ``Which is a Pascal's triangle number, sitting at the centre of row
% fourteen. The rook's question is a Pascal question, because the rook
% has two choices at every step and the recurrence has two parents.
% East or North. Two parents. Pascal.''
%  
% ``And the king?''
%  
% She smiled, and Mo recognised the smile from the Worpitzky episode,
% and knew he had asked the right question and was about to regret it.
%  
% ``The king is harder. He can step East, or North, or diagonally
% North-East. \emph{Three} options at every cell instead of two. So
% the recurrence has \emph{three} parents. Look. Let $D(i, j)$ be the
% number of king-paths to the cell $(i, j)$. Then,''
% \[
%   D(i, j) \;=\; D(i-1, j) \,+\, D(i, j-1) \,+\, D(i-1, j-1)
% \]
% ``because the king arrived from the west, from the south, or from
% the south-west diagonal. Three places he could have just come from.
% Sum them.''
%  
% She sketched the triangle quickly,
% \[
% \renewcommand{\arraystretch}{1.15}
% \setlength{\arraycolsep}{6pt}
% \begin{array}{cccccccc}
% 1 & 1 & 1 & 1 & 1 & 1 & 1 & 1 \\
% 1 & 3 & 5 & 7 & 9 & 11 & 13 & 15 \\
% 1 & 5 & 13 & 25 & 41 & 61 & 85 & 113 \\
% 1 & 7 & 25 & 63 & 129 & 231 & 377 & 575 \\
% 1 & 9 & 41 & 129 & 321 & 681 & 1289 & 2241 \\
% 1 & 11 & 61 & 231 & 681 & 1683 & 3653 & 7183 \\
% 1 & 13 & 85 & 377 & 1289 & 3653 & 8989 & 19825 \\
% 1 & 15 & 113 & 575 & 2241 & 7183 & 19825 & \mathbf{48639}
% \end{array}
% \]
%  
% ``Forty-eight thousand, six hundred and thirty-nine ways for the
% king to walk from corner to corner. Fourteen times more paths than
% the rook. The diagonal step gives him an enormous combinatorial
% inheritance.''
%  
% ``Does this triangle have a name?''
%  
% ``Delannoy. Henri Delannoy, late nineteenth century. He was a friend
% of Lucas, the same Lucas who gave us the Tower of Hanoi and the
% Fibonacci-cousin numbers. Lucas was the populariser. Delannoy was
% the quiet one who built the lattice that nobody noticed for fifty
% years.''
%  
% She tapped the diagonal of the table she had just drawn.
%  
% ``Now look at this spine, Mo, the bottom-left to top-right diagonal:
% one, three, thirteen, sixty-three, three twenty-one, sixteen
% eighty-three. The numbers along the centre. These are what people
% call the \emph{central Delannoy numbers}. They count the king-walks
% from corner to corner of an $n$ by $n$ board, and they are the
% heart of the whole structure. The forty-eight thousand we just
% computed is the eighth one. Notice that every term after the first
% is divisible by three. That is not a coincidence; it falls out of a
% small modular calculation, but never mind that tonight.''
%  
% ``Pascal has a spine too,'' Mo said cautiously, because he had been
% poked once before for forgetting.
%  
% ``Pascal has a spine,'' Eleanor agreed, pleased. ``The central
% column of Pascal's triangle, the binomial coefficients
% $\binom{2n}{n}$. One, two, six, twenty, seventy, two-fifty-two,
% nine-twenty-four. These count the rook-walks, the East-and-North
% paths only, with no diagonal. Forty-eight thousand for the king
% versus three thousand four hundred and thirty-two for the rook. The
% diagonal step is doing all that work.''
%  
% She drew the two spines on top of each other, the Delannoy above
% and the Pascal below.
% \[
% \renewcommand{\arraystretch}{1.25}
% \begin{array}{lcccccccc}
% \text{King } D(n,n) \!: & 1 & 3 & 13 & 63 & 321 & 1683 & 8989 & 48639 \\
% \text{Rook } \binom{2n}{n} \!: & 1 & 2 & 6 & 20 & 70 & 252 & 924 & 3432
% \end{array}
% \]
%  
% ``And here is the lovely bit. Each spine has a \emph{growth rate},
% the factor by which it multiplies as you move along it. Not the
% ratio of any two specific consecutive entries, mind you, but the
% limit those ratios crawl towards. For Pascal, you can compute
% $2/1, 6/2, 20/6, 70/20, 252/70, 924/252$, and you'll get two, three,
% three and a third, three and a half, three point six, three point
% six six, and so on. The ratios are climbing, but slowly, toward
% four. By the time you reach $3432/924$ you're at three point seven
% one, and the gap to four is still a quarter. The convergence is
% patient. But the destination is four, and four is just two squared,
% where the two comes from the rook having two choices at every step.
% Pascal's spine grows like $2^2$ per generation in the limit. Rational. The number you might guess.''
%  
% ``And Delannoy?''
%  
% ``Delannoy's spine grows like $(1 + \sqrt 2)^2$ in the limit. That is the silver mean, squared. Which is three plus two root two, which is
% roughly five point eight three. The actual ratios you see early
% on, three, four point three, four point eight, five point one, are
% crawling up toward that limit but they take a long time to get
% close. At $n$ equals seven you're at about five point four, still
% short. The pattern is the same as Pascal: ratios approach the
% asymptotic rate from below, never quite reaching it. But the
% {\em asymptote} itself is what we mean when we say `grows like'. The
% silver-mean-squared is what the spine \emph{wants} to grow by, and
% it gets closer and closer to that wish as $n$ increases.'' And with that she writes both rates side by side.
% \[
% \renewcommand{\arraystretch}{1.4}
% \begin{array}{lcl}
% \text{Pascal spine grows like} & 2^2 = 4 & \text{(rook, 3 choices)} \\
% \text{Delannoy spine grows like} & (1 + \sqrt 2)^2 = 3 + 2\sqrt 2 & \text{(king, 3 choices)}
% \end{array}
% \]
%  
% ``Both growth rates are \emph{squared expansion factors}, Mo.
% That's the pattern. For Pascal the expansion factor per step is
% just two, because the recurrence has two parents, and squaring it
% gives four. For Delannoy the expansion factor per step is the
% silver mean $1 + \sqrt 2$, because the recurrence has three parents
% and the third parent introduces a quadratic into the bookkeeping.
% The largest root of that quadratic is the silver mean, and squaring
% it gives the spine growth rate. The square is in there because to
% reach the spine entry at row $n$ you have to do $n$ worth of
% expansion on each of the two axes.''
%  
% ``Where does the silver mean come from?''
%  
% ``It comes from the characteristic equation of the recurrence.
% Suppose you ask Pascal's recurrence what number $\lambda$ makes
% $\lambda^n$ a solution. You get $\lambda = 2$, because two parents.
% Now ask Delannoy's recurrence the same question along its diagonal.
% You get $\lambda^2 = 2\lambda + 1$, which is the defining equation
% of the silver mean. The Fibonacci numbers come from
% $\lambda^2 = \lambda + 1$, and that gives you the golden mean.
% Delannoy is the silver-mean cousin of Fibonacci. The diagonal step
% in the lattice is what promotes the recurrence from `add the two
% parents' to `add the two parents and twice the running sum', and
% that quadratic difference is the difference between gold and
% silver. Different metal, same family.''
%  
% Mo thought about that for a while.
%  
% ``So the king walks faster than the rook because his shoes are made of
% silver.''
%  
% ``Silver instead of bronze, yes, in some metaphorical sense Mo
% which I shan't pursue.'' She paused. ``Although, while we're here:
% there's a bronze mean too. And copper. And nickel. The metallic
% means are an infinite family, each one the dominant root of
% $\lambda^2 = k\lambda + 1$ for $k = 1, 2, 3, \ldots$ Gold is $k = 1$,
% silver is $k = 2$, bronze is $k = 3$. Each one corresponds to a
% slightly different recurrence and a slightly different lattice
% problem. Pascal's lattice, in this sense, lives at $k = 0$, before
% the metals begin. Two choices, no diagonal, no metallic mean, just
% the integer two and its square.''
%  
% Mo looked at the triangle for a long moment. Then, because Mo was
% Mo, he asked the wrong question, which Eleanor had learnt long ago
% was the same as the right question.
%  
% ``Why is \emph{anyone} interested in counting paths a king can take?
% It's not like the king actually walks.''
%  
% >>>>>> end of annexed block <<<<<<
